\subsection{Genuine smooth principal-open affine lines}
\label{subsec:smooth-principal-open-lines}

\begin{theorem}[Actual polynomial-line smooth presentation]
\label{algebra:polynomial-line-standard-smooth}
\lean{EnvelopeProperty.polynomial_standardSmooth_one}
\leanok
For every commutative ring $R$, the actual polynomial algebra $R[T]$
is standard smooth of relative dimension one over $R$.
\end{theorem}
\begin{proof}
Present the one-variable polynomial algebra first as the multivariable
polynomial algebra on the actual one-element index set. Its true
polynomial generators have zero kernel, so the empty list of relations
spans that kernel. The empty differential is the actual identity on its
zero-dimensional module; its determinant and its image in the algebra
are one. Thus this is a genuine submersive presentation, with one
generator and no relations. Its true dimension is one. Transport the
presentation through the actual one-variable polynomial algebra
isomorphism. No assumed smooth-presentation or dimension certificate is used.
\end{proof}

\begin{theorem}[Nonzero domain elements have genuinely dense principal opens]
\label{algebra:domain-dense-principal-open}
\lean{EnvelopeProperty.dense_basicOpen_of_ne_zero}
\leanok
For every commutative domain $R$ and $0\ne r\in R$, the actual principal
open $D(r)\subseteq\operatorname{Spec}R$ is dense.
\end{theorem}
\begin{proof}
Every power of the nonzero element is a non-zero-divisor. The actual
localization map $R\to R[1/r]$ is consequently injective and has zero
kernel. Its kernel is therefore contained in the actual nilradical,
which is the genuine criterion for the induced spectrum map to have
dense range. The true spectrum image of this localization is exactly
$D(r)$. This proves density of the original principal open.
\end{proof}

\begin{theorem}[Actual principal-open line algebra]
\label{algebra:principal-open-line-ring}
\lean{EnvelopeProperty.AffineLine.OpenLineRing,
EnvelopeProperty.AffineLine.openLineEval,
EnvelopeProperty.AffineLine.openLineEval_algebraMap,
EnvelopeProperty.AffineLine.openLineRing_isDomain,
EnvelopeProperty.AffineLine.openLineRing_finitePresentation,
EnvelopeProperty.AffineLine.openLineRing_finiteType,
EnvelopeProperty.AffineLine.openLineRing_smooth,
EnvelopeProperty.AffineLine.openLineRing_isNoetherian,
EnvelopeProperty.AffineLine.openLineRing_standardSmooth_one,
EnvelopeProperty.AffineLine.openLineRing_isIntegrallyClosed,
EnvelopeProperty.AffineLine.openLineRing_isDedekindDomain}
\leanok
\uses{algebra:polynomial-line-standard-smooth}
Let $F$ be any field and $p\in F[T]$. The coordinate ring is the actual
localization $F[T,1/p]$. It is finitely presented and smooth over $F$,
and standard smooth of relative dimension one. It is noetherian.
If $p\ne0$, it is a domain, integrally closed and Dedekind.
Every parameter $t\in F$ with $p(t)\ne0$ gives an actual $F$-algebra
evaluator on this localization, agreeing with $f\mapsto f(t)$ on every
original polynomial.
\end{theorem}
\begin{proof}
The actual evaluation map sends $p$ to a unit of $F$, so the localization
universal property constructs the entire evaluator and proves its
restriction formula. Localization away from one element is finitely
presented, and the polynomial algebra is finitely presented over $F$;
composition gives finite presentation and finite type over $F$.
Finite type over a field gives noetherianity.

The preceding actual polynomial presentation has relative dimension
one; the actual principal localization is standard smooth of relative
dimension zero over the polynomial algebra. Composition of these genuine
presentations has relative dimension $1+0=1$, and also supplies smoothness.
For $p\ne0$, all its powers are non-zero-divisors, so genuine localization
domain inheritance applies. The polynomial ring over a field is integrally
closed and Dedekind; their actual localization theorems retain both
properties. No valuation, completeness, algebraic closure, characteristic
restriction or assumed curve witness is involved. The zero-polynomial
localization is allowed in the unconditional smoothness statements;
integrality and the normal-domain assertions explicitly require $p\ne0$.
\end{proof}

\begin{theorem}[Whole principal-open line and rational-point scheme maps]
\label{scheme:principal-open-line-map}
\lean{EnvelopeProperty.AffineLine.principalLineHom,
EnvelopeProperty.AffineLine.principalLineHom_algebraMap,
EnvelopeProperty.AffineLine.principalLineHom_surjective,
EnvelopeProperty.AffineLine.openLineEval_comp_principalLine,
EnvelopeProperty.AffineLine.openLineScheme,
EnvelopeProperty.AffineLine.openLineInclusion,
EnvelopeProperty.AffineLine.openLineInclusion_isOpenImmersion,
EnvelopeProperty.AffineLine.openLineInclusion_range,
EnvelopeProperty.AffineLine.openLineStructure,
EnvelopeProperty.AffineLine.openLineStructure_smooth_one,
EnvelopeProperty.AffineLine.openLineSchemePoint,
EnvelopeProperty.AffineLine.openLineSchemePoint_inclusion,
EnvelopeProperty.AffineLine.principalLineSchemeMap,
EnvelopeProperty.AffineLine.principalLineSchemeMap_isClosedImmersion,
EnvelopeProperty.AffineLine.principalLineSchemeMap_square,
EnvelopeProperty.AffineLine.principalLineSchemeMap_preimage_basicOpen,
EnvelopeProperty.AffineLine.openLineScheme_isIntegral,
EnvelopeProperty.AffineLine.ambientPrincipalPoint,
EnvelopeProperty.AffineLine.openLineSchemePoint_principalLine,
EnvelopeProperty.AffineLine.principalLineSchemeMap_zero}
\leanok
\uses{algebra:principal-open-line-ring,algebra:affine-line-substitution,scheme:affine-line-closed-immersion}
For any coordinate set, rational point $a$, direction $b$ and ambient
polynomial $g$, let $\alpha_b$ be the actual substitution
$X_i\mapsto a_i+b_iT$ and put $p=\alpha_b(g)$.
There is a genuine localized $F$-algebra map and its scheme map
\[
 F[X_i,1/g]\longrightarrow F[T,1/p],\qquad
 C_b=\operatorname{Spec}F[T,1/p]\longrightarrow D(g).
\]
For $b\ne0$ the whole map of coordinate rings is surjective, so this is
an actual closed immersion. The true structure morphism $C_b\to\operatorname{Spec}F$
is smooth of relative dimension one. If $p\ne0$, $C_b$ is an integral
scheme. Its actual open inclusion into the affine line has range $D(p)$.
Every principal-open preimage and every rational-point incidence is
computed on the entire original coordinate algebra. In particular,
$g(a)\ne0$ gives the whole original $D(g)$ rational point at parameter zero.
\end{theorem}
\begin{proof}
Localize the genuine whole substitution at $g$ and its actual image $p$.
The true localization map formula gives compatibility on every original
polynomial, and localized surjectivity follows from actual surjectivity
of the original nonconstant line substitution. Taking actual affine
spectra therefore gives the closed immersion and the whole commuting
square with the original affine-line immersion and both principal-open
inclusions. The localization spectrum theorem gives the actual open
immersion and its range $D(p)$.

The previous genuine ring smoothness statement is standard smooth of
relative dimension one. Passing through the true locally-standard-smooth
ring-hom property and its affine spectrum criterion gives the original
scheme structure morphism that same relative dimension. The actual
nonzero-localization domain makes its affine spectrum integral.

Given a parameter with $g(a+tb)\ne0$, both compositions of the actual
localized evaluators agree on every original polynomial by the full
line evaluation identity. The localization homomorphism uniqueness
theorem identifies the entire homomorphisms, so their genuine spectrum
maps commute on the entire rational point. At zero, $a+0b=a$ identifies
that point with the original one. Likewise, for any ambient polynomial
$h$, a prime in $C_b$ avoids the pullback of the actual element $h$ in
$F[X_i,1/g]$ exactly when it avoids the actual image of $\alpha_b(h)$.
This proves the whole good-principal-open preimage formula. No
coordinate-norm-only equality or assumed scheme incidence is used.
\end{proof}

\begin{theorem}[A genuine smooth open line with nonempty dense good restriction]
\label{scheme:smooth-principal-open-line-avoiding-function}
\lean{EnvelopeProperty.AffineLine.exists_parameter_avoiding,
EnvelopeProperty.AffineLine.exists_principal_open_line,
EnvelopeProperty.AffineLine.exists_principal_open_line_localized,
EnvelopeProperty.AffineLine.dense_openLine_basicOpen,
EnvelopeProperty.AffineLine.nonempty_openLine_basicOpen,
EnvelopeProperty.AffineLine.exists_smooth_principal_open_line,
EnvelopeProperty.AffineLine.exists_smooth_principal_open_line_localized}
\leanok
\uses{algebra:affine-line-avoids-polynomial,algebra:domain-dense-principal-open,scheme:principal-open-line-map}
Let $F$ be an infinite field, with a nonempty coordinate set, and let
$a\in D(g)(F)$ be an actual rational point. For every actual nonzero
function $f\in F[X_i,1/g]$, there is a genuine nonzero direction $b$
such that the preceding $C_b$ is an integral smooth affine scheme of
relative dimension one, closed in $D(g)$, containing the entire original
point $a$. The actual scheme preimage of $D(f)$ in $C_b$ is nonempty and
dense. The polynomial-function specialization also provides an actual
rational parameter outside the zeros of both restricted polynomials.
\end{theorem}
\begin{proof}
First let the good function be a nonzero ambient polynomial $h$.
The actual line-avoidance construction supplies a nonzero direction
$b$ with $\alpha_b(h)\ne0$. Since $g(a)\ne0$ and evaluation at zero
is the entire original evaluation, $\alpha_b(g)\ne0$ as well.
Their product is a nonzero polynomial. Over the infinite field it cannot
vanish at every parameter, yielding a genuine rational parameter where
both restrictions are nonzero. Its actual localized evaluator gives a
prime in the good open, proving its nonemptiness. Injectivity of the
nonzero principal localization preserves the nonzero good numerator,
so the domain density theorem also makes this good open dense.

For the actual localized function $f$, choose its genuine localization
representation
\[
 f\,g^n=q\quad\text{in }F[X_i,1/g].
\]
The actual image of $g^n$ is a unit, so $f\ne0$ forces the original
numerator $q\ne0$. Apply the polynomial construction to $q$, retaining
$\alpha_b(g)\ne0$ and $\alpha_b(q)\ne0$. Apply the whole localized
line homomorphism to the preceding equality. If its image of $f$ were
zero, its image of $q$ would be zero; injectivity of the actual domain
localization would then force $\alpha_b(q)=0$, a contradiction.
Thus the genuine localized line restriction of $f$ is nonzero.
Its principal open is dense by the actual domain theorem, and a dense
subset of the nonempty integral affine spectrum is nonempty.
The true scheme pullback on prime supports identifies this principal
open with the whole preimage of $D(f)$.

Finally use the actual closed immersion, smooth relative-dimension-one
structure morphism, integral-spectrum theorem and full zero-parameter
incidence proved above. These give the stated construction directly,
without an assumed line, curve, smoothness or nonempty-open certificate.
For FHX Proposition~5.22 this supplies the open-line construction after
the genuinely constructed complete algebraically closed field and its
rigid point are available. It does not yet supply normalization or
reduction of a finite preimage, dominance of every irreducible component,
finite flatness of the resulting normalized map, finite bad scheme-point
sets, or subharmonic curve potential theory. Those are distinct producers.
\end{proof}
